\begin{claim}{B.39}{Proved}{Complete genus-three projective-plane quotient catalogue}
There are exactly $11$ conjugacy classes of mixed finite genus-three
actions whose positive quotient is a sphere and whose induced
involution on that sphere is fixed-point-free.  Their full quotients
are projective planes with cone points.  Four classes are identified
with B.34 and seven are additional.  Consequently B.37 and the present
claim give exactly $156$ mixed classes with spherical positive quotient.
\end{claim}
\begin{proof}
A fixed-point-free sphere involution has projective-plane quotient:
its quotient is a closed surface of Euler characteristic one, and
its orientation cover is the given sphere.  Cone points are exchanged
in pairs.  Thus each multiplicity in the positive cone multiset must
be even.  Inspection of the complete input in B.33 leaves exactly
$\mathrm O3.4$, $\mathrm O3.9$, $\mathrm O3.11$, $\mathrm O3.15$,
$\mathrm O3.21$, $\mathrm O3.22$, and $\mathrm O3.25$.
Their paired cone orders are respectively
$(2,2,2,2)$, $(2,2,2)$, $(4,4)$, $(2,6)$, $(2,4)$, $(2,4)$,
and $(2,3)$.

For a projective plane with cone orders $n_1,\ldots,n_l$ use
\[
 \Lambda=\langle x_1,\ldots,x_l,d\mid
       x_i^{n_i}=1,\quad x_1\cdots x_l d^2=1\rangle .   \tag{B.39.1}
\]
Set $X_i=x_i$, $Y_i=d x_i d^{-1}$ and
$A=d^2=(X_1\cdots X_l)^{-1}$.  Its orientation double has the
spherical vector
\[
        v=(X_1,\ldots,X_l,Y_l^{-1},\ldots,Y_1^{-1}),
                \qquad\prod_iX_i=\prod_iY_i.            \tag{B.39.2}
\]
Conjugation by $d$ is
\[
       X_i\longmapsto Y_i,\qquad
       Y_i\longmapsto A X_iA^{-1}.                     \tag{B.39.3}
\]
For every published $\operatorname{Aut}(H)$-normalized vector in
these seven cases, the program tests that (B.39.3) descends to an
automorphism $\alpha$ of its finite positive image $H$.  The same
Cayley-graph propagation as in B.37 gives an exact kernel-invariance
test.  It also verifies
$\alpha^2=\operatorname{Inn}(A)$ and $\alpha(A)=A$.
The full finite image is then the set $H\times\{0,1\}$ with product
\[
 (h,\epsilon)(k,\delta)=
       \bigl(h\alpha^\epsilon(k) A^{\epsilon\delta},
                    \epsilon+\delta\bmod2\bigr).      \tag{B.39.4}
\]
The two identities just verified prove associativity; the element
$d=(1,1)$ has square $A$ and induces $\alpha$ on $H$.
Thus (B.39.4) is exactly $\Lambda/\ker\theta$, where
$\theta$ is the positive epimorphism (B.39.2).

For clarity, the full source fusion is as follows.  The colored
mapping class group of the punctured projective plane is generated
by the colored disc braids and by sliding each puncture once around
a one-sided loop.  One direct proof cuts the projective plane along
a projective line disjoint from the marked disc: a moving puncture
crossing the cut is separated off as a one-sided slide, after which
all the remaining tracks form a braid in the disc.  Equivalently,
forgetting punctures successively expresses their point-pushing
loops by peripheral loops and the one-sided generator.  The former
are the pure disc braids.  The unmarked projective-plane mapping
class group is trivial: isotope the image of a projective line to
that line, adjust its parametrization by a projective rotation, and
apply the disc Alexander isotopy after cutting.  This proves that
these point motions exhaust all mapping classes.  Color-preserving
permutations are supplied by the same equal-color braid generators
used in B.37.

The slide of the last puncture has the exact substitutions
\[
 x_i\longmapsto x_i\quad(i<l),\qquad
 x_l\longmapsto(x_l d)x_l^{-1}(x_l d)^{-1},\qquad
 d\longmapsto x_l d.                                   \tag{B.39.5}
\]
They are obtained by tracking the meridian through a neighborhood
of that puncture and the one-sided loop.  The two-cone case, including
its geometric puncture-slide interpretation, is verified independently
in \cite[Lemma 3.19]{C_GromadzkiSzepietowski2016}.  Other punctures
outside this neighborhood are fixed, so the same formula applies
to all $l$ used here.  It respects the product relation directly:
if $B=x_1\cdots x_{l-1}$ and $x=x_l$, then
\[
 B(x d x^{-1}d^{-1}x^{-1})(x d)^2=B x d^2=1.
\]
In the finite positive group its effect is
\[
 x_l\longmapsto x_l\alpha(x_l)^{-1}x_l^{-1},\qquad
 \alpha\longmapsto\operatorname{Inn}(x_l)\alpha,
 \qquad A\longmapsto x_l\alpha(x_l)A.                  \tag{B.39.6}
\]
To slide any other puncture, a specified sequence of adjacent braids
moves it to the last position, applies (B.39.6), and moves it back.
The full word preserves its color even when the intermediate braids
exchange different colors.  Equations (B.39.5)--(B.39.6) and the disc
braids therefore give exactly the source mapping-class action.
The kernel-to-cover conjugacy argument of B.37 applies without change.

Executing these finite orbits gives respectively $1,3,3,1,2,0,1$
classes in the seven signature/group pairs above.  For every class,
\texttt{catalogue\_genus3\_projective\_quotients.py} constructs the
regular representation of (B.39.4), checks every relation in
(B.39.1), injectivity of all local cyclic groups, and Euler
characteristic $-4$.  The finite point-coset calculation gives every
partial section and its multiplier, centralizer and extensions.
No negative element has a fixed point, since its image on the
positive quotient sphere is fixed-point-free.  Thus all nontrivial
sections in these rows have positive cyclic image and are point
sections; their surface-group centralizers are trivial.

Finally the four projective-plane groups in B.34 are compared using
their actual cone generators and conjugation by their crosscap
element.  After a fixed ordering braid, their doubled vectors are
identified in the finite kernel orbits by the rooted colored Cayley
graph isomorphism test.  The resulting exact identifications are
\[
 \mathrm{G3H}_{002}\leftrightarrow\mathrm{G3P}_{001},\quad
 \mathrm{G3H}_{011}\leftrightarrow\mathrm{G3P}_{002},\quad
 \mathrm{G3H}_{029}\leftrightarrow\mathrm{G3P}_{008},\quad
 \mathrm{G3H}_{039}\leftrightarrow\mathrm{G3P}_{010}.
                                                        \tag{B.39.7}
\]
All seven applicable section-count columns agree under this
bijection.  Hence precisely seven classes are additional to B.34.
The quotient involution on the positive sphere has a fixed circle
in B.37 and has empty fixed set here; a conjugacy induces a sphere
conjugacy carrying these fixed sets to each other.  Thus none of
these eleven rows can be conjugate to a B.37 row, proving the total
$145+11=156$.
\end{proof}

Here $p$ counts nontrivial partial section classes, and $m,o$ are
$|\Phi|$ and $|K\backslash\Phi|$.  The final column $i$ counts negative
involutions in the displayed finite group; all of them act freely.
\[
\begin{array}{c|r|l|r|r|r|r}
 \text{class}&|K|&\text{cone orders}&p&m&o&i\\\hline
 001&4&(2,2,2,2)&8&8&4&2\\
 002&8&(2,2,2)&12&12&3&0\\
 003&8&(2,2,2)&12&12&3&2\\
 004&8&(2,2,2)&12&12&3&4\\
 005&8&(4,4)&8&4&2&0\\
 006&8&(4,4)&8&4&2&2\\
 007&8&(4,4)&8&4&2&4\\
 008&12&(2,6)&12&8&2&2\\
 009&16&(2,4)&16&12&2&0\\
 010&16&(2,4)&16&12&2&0\\
 011&24&(2,3)&20&20&2&4
\end{array}                                             \tag{B.39.8}
\]
Every row has zero full splittings, since its full image contains a
negative element and no such element fixes a point.  The JSON
\texttt{type\_3\_0\_projective\_quotient\_sections.json} includes all
finite permutations, local generators, multipliers and section
incidences behind this table.
